\begin{abstract}
Generative inpainting can remove or replace a small part of a photograph while leaving most of the scene intact. A forensic system must identify both whether an image was edited and where the edit occurred. We study a single-image detector that combines semantic features with high-pass residual features and uses a resized whole-image view alongside overlapping native-resolution tiles. The larger model combines ConvNeXt and frozen DINOv2 features with a residual encoder and cross-attention decoder. A separate training configuration assigns crop labels according to the presence of annotated edit pixels. In a saved evaluation on a previously inspected 100-image development subset, an ensemble obtained 70.0\% image accuracy, 0.7452 AUROC, and 34.89\% mean localization IoU; mean IoU was 8.94\% for edits below 2\% of the image area. A local 10K target scorecard lists 91.20\% accuracy and 47.30\% mean IoU, but its metrics conflict with its accompanying confusion counts and its prediction-level artifacts have not been located. We present those figures as unverified targets, not a validated head-to-head result. Our findings motivate source-grouped independent evaluation, consistent metric aggregation, and separate analysis of small edits, missed edits, and false alarms.
\end{abstract}
